\chapter{AES tabele}
\label{app:aes_tabele}

\section{S-boks}
\label{app:sbox}

Tabela \ref{tab:sbox} sadrži kompletan S-boks AES algoritma \cite{fips197}. Vrsta se bira gornjim, a kolona donjim polubajtom ulaznog bajta. Za ulaz $\{53\}$ vrednost se nalazi u vrsti \texttt{5} i koloni \texttt{3}, i iznosi $\{ED\}$, što je primer korišćen u odeljku o SubBytes transformaciji. Inverzni S-boks u ovom radu nije potreban, budući da GCM režim koristi AES isključivo u smeru šifrovanja.

\begin{table}[H]
\caption{AES S-boks: vrsta je gornji, kolona donji polubajt ulaza (heksadecimalno)}
\centering
\renewcommand{\arraystretch}{1.15}
\setlength{\tabcolsep}{4pt}
\resizebox{0.92\textwidth}{!}{
% Automatski generisano skriptom gen_tables.py -- AES S-box (FIPS 197)
\begin{tabular}{|c|c|c|c|c|c|c|c|c|c|c|c|c|c|c|c|c|}
\hline
\cellcolor{gray!30} & \cellcolor{gray!30}\texttt{0} & \cellcolor{gray!30}\texttt{1} & \cellcolor{gray!30}\texttt{2} & \cellcolor{gray!30}\texttt{3} & \cellcolor{gray!30}\texttt{4} & \cellcolor{gray!30}\texttt{5} & \cellcolor{gray!30}\texttt{6} & \cellcolor{gray!30}\texttt{7} & \cellcolor{gray!30}\texttt{8} & \cellcolor{gray!30}\texttt{9} & \cellcolor{gray!30}\texttt{A} & \cellcolor{gray!30}\texttt{B} & \cellcolor{gray!30}\texttt{C} & \cellcolor{gray!30}\texttt{D} & \cellcolor{gray!30}\texttt{E} & \cellcolor{gray!30}\texttt{F} \\ \hline
\cellcolor{gray!30}\texttt{0} & \texttt{63} & \texttt{7C} & \texttt{77} & \texttt{7B} & \texttt{F2} & \texttt{6B} & \texttt{6F} & \texttt{C5} & \texttt{30} & \texttt{01} & \texttt{67} & \texttt{2B} & \texttt{FE} & \texttt{D7} & \texttt{AB} & \texttt{76} \\ \hline
\cellcolor{gray!30}\texttt{1} & \texttt{CA} & \texttt{82} & \texttt{C9} & \texttt{7D} & \texttt{FA} & \texttt{59} & \texttt{47} & \texttt{F0} & \texttt{AD} & \texttt{D4} & \texttt{A2} & \texttt{AF} & \texttt{9C} & \texttt{A4} & \texttt{72} & \texttt{C0} \\ \hline
\cellcolor{gray!30}\texttt{2} & \texttt{B7} & \texttt{FD} & \texttt{93} & \texttt{26} & \texttt{36} & \texttt{3F} & \texttt{F7} & \texttt{CC} & \texttt{34} & \texttt{A5} & \texttt{E5} & \texttt{F1} & \texttt{71} & \texttt{D8} & \texttt{31} & \texttt{15} \\ \hline
\cellcolor{gray!30}\texttt{3} & \texttt{04} & \texttt{C7} & \texttt{23} & \texttt{C3} & \texttt{18} & \texttt{96} & \texttt{05} & \texttt{9A} & \texttt{07} & \texttt{12} & \texttt{80} & \texttt{E2} & \texttt{EB} & \texttt{27} & \texttt{B2} & \texttt{75} \\ \hline
\cellcolor{gray!30}\texttt{4} & \texttt{09} & \texttt{83} & \texttt{2C} & \texttt{1A} & \texttt{1B} & \texttt{6E} & \texttt{5A} & \texttt{A0} & \texttt{52} & \texttt{3B} & \texttt{D6} & \texttt{B3} & \texttt{29} & \texttt{E3} & \texttt{2F} & \texttt{84} \\ \hline
\cellcolor{gray!30}\texttt{5} & \texttt{53} & \texttt{D1} & \texttt{00} & \texttt{ED} & \texttt{20} & \texttt{FC} & \texttt{B1} & \texttt{5B} & \texttt{6A} & \texttt{CB} & \texttt{BE} & \texttt{39} & \texttt{4A} & \texttt{4C} & \texttt{58} & \texttt{CF} \\ \hline
\cellcolor{gray!30}\texttt{6} & \texttt{D0} & \texttt{EF} & \texttt{AA} & \texttt{FB} & \texttt{43} & \texttt{4D} & \texttt{33} & \texttt{85} & \texttt{45} & \texttt{F9} & \texttt{02} & \texttt{7F} & \texttt{50} & \texttt{3C} & \texttt{9F} & \texttt{A8} \\ \hline
\cellcolor{gray!30}\texttt{7} & \texttt{51} & \texttt{A3} & \texttt{40} & \texttt{8F} & \texttt{92} & \texttt{9D} & \texttt{38} & \texttt{F5} & \texttt{BC} & \texttt{B6} & \texttt{DA} & \texttt{21} & \texttt{10} & \texttt{FF} & \texttt{F3} & \texttt{D2} \\ \hline
\cellcolor{gray!30}\texttt{8} & \texttt{CD} & \texttt{0C} & \texttt{13} & \texttt{EC} & \texttt{5F} & \texttt{97} & \texttt{44} & \texttt{17} & \texttt{C4} & \texttt{A7} & \texttt{7E} & \texttt{3D} & \texttt{64} & \texttt{5D} & \texttt{19} & \texttt{73} \\ \hline
\cellcolor{gray!30}\texttt{9} & \texttt{60} & \texttt{81} & \texttt{4F} & \texttt{DC} & \texttt{22} & \texttt{2A} & \texttt{90} & \texttt{88} & \texttt{46} & \texttt{EE} & \texttt{B8} & \texttt{14} & \texttt{DE} & \texttt{5E} & \texttt{0B} & \texttt{DB} \\ \hline
\cellcolor{gray!30}\texttt{A} & \texttt{E0} & \texttt{32} & \texttt{3A} & \texttt{0A} & \texttt{49} & \texttt{06} & \texttt{24} & \texttt{5C} & \texttt{C2} & \texttt{D3} & \texttt{AC} & \texttt{62} & \texttt{91} & \texttt{95} & \texttt{E4} & \texttt{79} \\ \hline
\cellcolor{gray!30}\texttt{B} & \texttt{E7} & \texttt{C8} & \texttt{37} & \texttt{6D} & \texttt{8D} & \texttt{D5} & \texttt{4E} & \texttt{A9} & \texttt{6C} & \texttt{56} & \texttt{F4} & \texttt{EA} & \texttt{65} & \texttt{7A} & \texttt{AE} & \texttt{08} \\ \hline
\cellcolor{gray!30}\texttt{C} & \texttt{BA} & \texttt{78} & \texttt{25} & \texttt{2E} & \texttt{1C} & \texttt{A6} & \texttt{B4} & \texttt{C6} & \texttt{E8} & \texttt{DD} & \texttt{74} & \texttt{1F} & \texttt{4B} & \texttt{BD} & \texttt{8B} & \texttt{8A} \\ \hline
\cellcolor{gray!30}\texttt{D} & \texttt{70} & \texttt{3E} & \texttt{B5} & \texttt{66} & \texttt{48} & \texttt{03} & \texttt{F6} & \texttt{0E} & \texttt{61} & \texttt{35} & \texttt{57} & \texttt{B9} & \texttt{86} & \texttt{C1} & \texttt{1D} & \texttt{9E} \\ \hline
\cellcolor{gray!30}\texttt{E} & \texttt{E1} & \texttt{F8} & \texttt{98} & \texttt{11} & \texttt{69} & \texttt{D9} & \texttt{8E} & \texttt{94} & \texttt{9B} & \texttt{1E} & \texttt{87} & \texttt{E9} & \texttt{CE} & \texttt{55} & \texttt{28} & \texttt{DF} \\ \hline
\cellcolor{gray!30}\texttt{F} & \texttt{8C} & \texttt{A1} & \texttt{89} & \texttt{0D} & \texttt{BF} & \texttt{E6} & \texttt{42} & \texttt{68} & \texttt{41} & \texttt{99} & \texttt{2D} & \texttt{0F} & \texttt{B0} & \texttt{54} & \texttt{BB} & \texttt{16} \\ \hline
\end{tabular}

}
\label{tab:sbox}
\end{table}

\section{Konstante rundi ekspanzije ključa}
\label{app:rcon}

Tabela \ref{tab:rcon} sadrži konstante $Rcon_i = (x^{i-1}, \{00\}, \{00\}, \{00\})$ koje se koriste u ekspanziji ključa. Prikazan je prvi bajt $x^{i-1}$, izračunat u polju $\mathrm{GF}(2^8)$. AES-256 koristi konstante do $Rcon_7$, dok AES-128 koristi svih deset.

\begin{table}[H]
\caption{Rcon konstante ekspanzije ključa (prvi bajt, heksadecimalno)}
\centering
\renewcommand{\arraystretch}{1.15}
\resizebox{0.7\textwidth}{!}{
% Automatski generisano -- Rcon konstante
\begin{tabular}{|c|c|c|c|c|c|c|c|c|c|c|}
\hline
\cellcolor{gray!30}$i$ & \cellcolor{gray!30}1 & \cellcolor{gray!30}2 & \cellcolor{gray!30}3 & \cellcolor{gray!30}4 & \cellcolor{gray!30}5 & \cellcolor{gray!30}6 & \cellcolor{gray!30}7 & \cellcolor{gray!30}8 & \cellcolor{gray!30}9 & \cellcolor{gray!30}10 \\ \hline
\cellcolor{gray!30}$Rcon_i$ & \texttt{01} & \texttt{02} & \texttt{04} & \texttt{08} & \texttt{10} & \texttt{20} & \texttt{40} & \texttt{80} & \texttt{1B} & \texttt{36} \\ \hline
\end{tabular}

}
\label{tab:rcon}
\end{table}
